\section{Semantic preservation: model, proof and implementation obligations}\label{sec:semantics}
We give a preservation argument for a restricted logical model, then identify
what is additionally required of the implemented compiler and runtime. This
separates a mathematical property of Cell reuse from a claim that the current
Rust implementation implements all of egglog correctly. The latter is not
established by this paper or by its timing experiments.

\subsection{Reference model and observations}
Fix a typed program with constructor declarations, positive conjunctive queries,
equality tests, and deterministic partial scalar primitives. RHS instructions
construct terms or request equality. Primitive failure has the same meaning in
both interpreters. Partial failures are admitted as query rejection; RHS primitive
evaluations must succeed for every admitted binding. Aborting execution with
partially committed effects is outside the theorem. The fragment excludes deletion, subsumption, mutable function
updates, identity-dependent operations, observable extraction tie choices and
external effects. A primitive whose result depends on interning history is not
pure in this model. Initial facts and the original source schedule are fixed.
Each Run reads its boundary snapshot, obtains all distinct rule/input bindings,
and applies their RHS instructions; new matches become available at the next
source-scheduled activation. This is the reference model used in the theorem,
not a definition asserted to cover every native egglog scheduler or command.

A logical state is a finite typed presentation of constructor facts and explicit
equalities. Its completed state uses the least equivalence relation containing
the explicit equalities and closed under constructor congruence. Concrete IDs
are names, not observations. Two completed states are related, written
$S\simeq T$, when a sort-preserving bijection of their equality classes preserves
all constructor rows and distinguished initial values, while scalar values
have the same primitive-observable meaning. Thus equality of row counts or
extraction costs alone does not establish $\simeq$.

Let $\alpha(C)$ interpret a Cell runtime state: decode each live fact through its
template and actual instance slots, resolve ports through the current equality
relation, and project away compiler-private Entry and cursor data. At an
observation boundary, no unresolved repair may affect the projected relations.
For the restricted fragment, $\alpha$ contains every logical fact; retained
private storage is not itself a logical addition.

\subsection{Three local preservation lemmas}
\paragraph{Lemma 1 (query-prefix reuse).}
Let a query be $Q=A_1\land\cdots\land A_k$. Suppose an Entry supplies a partial
binding $\theta$ and valid witnesses proving $A_1,\ldots,A_j$ in the same snapshot
used by the remaining query. Evaluating the remaining atoms with $\theta$ yields
exactly those full bindings of $Q$ extending $\theta$.

\emph{Proof.} A returned extension satisfies the remaining atoms by evaluation
and the prefix by its witnesses, so it satisfies $Q$. Conversely, restrict any
full satisfying binding extending $\theta$ to the operands of each remaining
atom. The ordinary exhaustive evaluator admits that extension at every step.
This proves equality for this entry, not completeness of the collection of
entries. To replace enumeration, the union of specialized and residual routes
must cover every reference binding, with the prescribed duplicate ownership.
In particular, a binding enabled by a late input or a new equality cannot be
omitted merely because no historical cursor reaches it. $\square$

\paragraph{Lemma 2 (instruction and witness reuse).}
Suppose corresponding operand slots denote the same logical values. A constructor
instruction $F(a_1,\ldots,a_n)$ may reuse an actual result witness $u$ if the
current state proves $F(a_1,\ldots,a_n)=u$. Replacing the instruction by this
witness leaves the completed logical state and its result unchanged. Copying an
existing operand and evaluating an identical pure scalar primitive also preserve
the relation.

\emph{Proof.} Inserting a constructor row already represented at that key adds
no new logical fact after congruence closure; its result is the same equality
class. A physical allocation or producer identity is not observable in this
fragment. Copies preserve their operands and pure primitives have identical
partial interpretations. These arguments apply instruction by instruction.
The witness must name the actual result, which may belong to another Cell;
a producer/op address without that evidence does not satisfy the premise.
An explicit union can be omitted only if its operands are already provably
equal. Constructing a previously absent term is not justified by changing a
Cell's state label. $\square$

\paragraph{Lemma 3 (local congruence completion).}
For a finite set of constructor rows, canonical-key rebuilding and a complete
fair evaluation of the implications
\[
 F(\mathbf a)=u\;\land\;F(\mathbf b)=v\;\land\;
 \bigwedge_{i=1}^{n} a_i\equiv b_i\quad\Longrightarrow\quad u\equiv v
\]
produce the same least congruence extending the explicit unions.

\emph{Proof.} Every implication adds an equality required by congruence, so
no step exceeds the least closed relation. Both procedures monotonically merge
classes in a finite carrier. At their fixed points all applicable implications
hold; any such closed relation contains the least one. Soundness and closure
therefore give equality of the results. Fairness means every persistently
applicable implication is eventually examined. A template may execute a static
subset of these implications, provided the residual mechanism covers all other
pairs, including pairs newly enabled across Cell boundaries. Equality of only
one argument of a multi-argument constructor is insufficient. A digest or Bloom
filter is not a proof of the conjunctive premise. $\square$

\subsection{Boundary simulation theorem}
\paragraph{Theorem (restricted-fragment preservation).}
Consider the reference interpreter above and a Cell interpreter for the same
program. Assume (i) their initial projections are related; (ii) the typed lowering
preserves each query atom, RHS instruction and source control boundary;
(iii) specialized entries plus residual enumeration yield exactly the reference
rule/input bindings; (iv) reuse satisfies Lemmas 1 and 2; and (v) completion
satisfies Lemma 3 before each observation. Then the two interpreters have related
completed states at every corresponding scheduled boundary. Source checks over
the modeled relations have identical outcomes.

\emph{Proof.} Induct on source-scheduled activations. Initially (i) establishes
the relation. At an induction boundary, the class bijection maps each query
binding to a corresponding binding; positive atoms and pure guards agree.
Conditions (ii)--(iii) and Lemma 1 supply the same distinct rule/input pairs.
Apply corresponding RHS instructions to these pairs. Fresh constructor results
extend the correspondence with their logical constructor terms, already existing
results preserve it, and union requests add the same logical equalities.
Lemma 2 permits justified reuse steps to stutter. In this fragment the resulting
set of facts and unions is independent of the enumeration order. Lemma 3 gives
the same least completed presentation, extending the class correspondence through
any merges. Publication exposes exactly that presentation at the next boundary,
so the induction hypothesis is restored. A check is a query over related states
and therefore agrees. $\square$

For a fixed sequence or fixed repeat count this argument supplies the complete
finite trace. For conditionals, until, or saturation it additionally requires
agreement of the source control predicate, including any change/progress flag
used for stopping. Physical allocation counts, cursor advancement and the
number of cached hits cannot substitute for that predicate. No termination claim
is made for unbounded programs; if matching control choices hold, every finite
prefix has the same observations.

\subsection{Composition and the obligations beyond the fragment}
A compiled dominated-to-dominator transition is a finite sequence of the local
steps above. Induction on that sequence proves preservation when each step's
premises hold and original Run boundaries remain in place. Compatible slot
layouts change representation rather than logical values. Multiple derivations
can share a destination only with their own entry premises and result mapping
preserved. Structural inclusion alone does not establish those premises or
justify additional effects. Likewise, query workers may reorder evaluation of
a fixed snapshot without changing its complete answer set; moving writes into
workers requires a separate commutation or serialization argument.

For the full supported executable language, use a stronger simulation relation
that also preserves visible versus subsumed rows, mutable function contents,
primitive-observable representation state, and future query/delta eligibility.
In that setting, equality of answer sets is insufficient. For example, two writes
with a last-writer merge can yield different results under a different binding
order. Deleting a witnessed row invalidates reuse even if its old slot remains
allocated. A cached mutable read can become wrong without any input equality
changing. The required extension preserves the ordered observable effect stream,
including merge evaluation, failure behavior and extraction choices, or proves
that each reordered pair commutes. The monotone theorem does not provide that
extension automatically.

\begin{table}[ht]\centering\small
\begin{tabularx}{\linewidth}{lXX}\toprule
Proof boundary & Required implementation obligation & Current evidence and limit\\\midrule
Source lowering & Every admitted source command has matching reference semantics; private marker views expand to the source query & Typed admission and differential tests; no complete compiler-correctness proof\\
Entry coverage & All reference bindings, including union-enabled and late-input matches, survive partition and deduplication & Residual routes and Run audits; not a proof for every compiled plan\\
Effects and repair & Actual witnesses, complete congruence and correct mutable/deletion behavior & Runtime checks and selected tests; Lemma 3 proves the abstract complete closure, not all repair code\\
Source control & Matching visibility, delta ownership, stopping and effect order & Selected boundary audits; final-state tests alone do not discharge this obligation\\
Bake composition & Check each inherited premise, result mapping and boundary & Typed catalog checks and observed replay; no blanket soundness claim for all generated catalogs\\\bottomrule
\end{tabularx}
\caption{The abstract preservation argument exposes obligations that must be
proved or validated in the implementation; their statement is not their discharge.}
\end{table}

The current upstream final-state passes and the source-specific sum-bits
acceptance contract provide evidence at particular inputs. They neither prove
all five theorem premises for the implementation nor establish equivalence to
every native version. Conversely, a different physical representative or a
separately validated extraction tie is not by itself a semantic counterexample.
A future unconditional compiler-correctness claim needs a formal native reference
semantics for the admitted fragment and proofs connecting the actual lowering,
matching, effects, completion and source control to it.
